\begin{remark*}[Exposition]
A metric is the distance function itself.  A metric space is the set together
with that function.  This section records standard metric functions and then
states, one by one, that each function satisfies the metric axioms.
\end{remark*}

\subsection{Norms Induce Metrics}

\begin{definitionbox}{Definition --- Norm}
\begin{definition}[Norm]\label{def:norm-on-real-vector-space}
Let \(V\) be a real vector space.  A \emph{norm} on \(V\) is a function
\[
  \|\cdot\|:V\to\mathbb{R}_{\geq 0}
\]
such that for all \(x,y\in V\) and all \(\lambda\in\mathbb{R}\),
\begin{enumerate}
  \item \(\|x\|=0\) if and only if \(x=0\)
        \hfill\textit{(positive definiteness)};
  \item \(\|\lambda x\|=|\lambda|\,\|x\|\)
        \hfill\textit{(absolute homogeneity)};
  \item \(\|x+y\|\leq \|x\|+\|y\|\)
        \hfill\textit{(triangle inequality)}.
\end{enumerate}
\end{definition}
\end{definitionbox}

\begin{remark*}[Standard quantified statement]
\[
\begin{aligned}
\operatorname{IsNorm}(V,\|\cdot\|)
\Longleftrightarrow\;
&\|\cdot\|:V\to\mathbb{R}_{\geq 0}
\\
&{}\land
\forall x\in V{:}\; \|x\|=0 \Longleftrightarrow x=0
\\
&{}\land
\forall \lambda\in\mathbb{R}\;\forall x\in V{:}\;
\|\lambda x\|=|\lambda|\,\|x\|
\\
&{}\land
\forall x,y\in V{:}\; \|x+y\|\leq \|x\|+\|y\|.
\end{aligned}
\]
\end{remark*}

\begin{remark*}[Predicate reading]
\[
\operatorname{IsNorm}(V,\|\cdot\|)
\]
means that \(\|\cdot\|\) is a norm on the real vector space \(V\).
\end{remark*}

\begin{remark*}[Interpretation]
A norm measures the size of one vector.  The metric induced by a norm measures
the size of the displacement vector \(x-y\).  Thus a norm is a length
structure, and the induced metric is the corresponding distance structure.
\end{remark*}

\DefinitionalRoot

\begin{theorembox}{Theorem}
\begin{theorem}[A Norm Induces a Metric]\label{thm:norm-induces-metric}
Let \(V\) be a real vector space and let \(\|\cdot\|\) be a norm on \(V\).
Define
\[
  d_{\|\cdot\|}:V\times V\to\mathbb{R}_{\geq 0},
  \qquad
  d_{\|\cdot\|}(x,y):=\|x-y\|.
\]
Then \(d_{\|\cdot\|}\) is a metric on \(V\).

\smallskip
\noindent
\hyperref[prf:norm-induces-metric]{\textit{Go to proof.}}
\end{theorem}
\end{theorembox}

\begin{remark*}[Standard quantified statement]
\[
\operatorname{IsNorm}(V,\|\cdot\|)
\Longrightarrow
\operatorname{IsMetric}(V,d_{\|\cdot\|}),
\qquad
d_{\|\cdot\|}(x,y)=\|x-y\|.
\]
\end{remark*}

\begin{remark*}[Predicate reading]
\[
\operatorname{IsNorm}(V,\|\cdot\|)
\Longrightarrow
\operatorname{IsMetric}(V,d_{\|\cdot\|})
\]
means that a norm on \(V\) induces a metric on \(V\) by distance through
vector differences.
\end{remark*}

\begin{remark*}[Interpretation]
The metric axioms are exactly the norm axioms applied to displacement
vectors.  Zero distance says the displacement \(x-y\) has zero norm, hence is
the zero vector.  Symmetry follows from absolute homogeneity with scalar
\(-1\).  The triangle inequality follows by writing \(x-y=(x-z)+(z-y)\).
\end{remark*}

\begin{dependencies}
\begin{itemize}
  \item \hyperref[def:norm-on-real-vector-space]{Definition: Norm}
  \item \hyperref[def:metric-space]{Definition: Metric Space}
\end{itemize}
\end{dependencies}
